\section{A population-mismatch illustration}
\label{sec:eca}

Changing the reference can also change the population being assessed.
A constructed external-control example makes this distinction explicit.
A fixed synthetic registry supplies a single-arm trial and three controls
per enrollee, with enrollment favoring heavily pretreated strata and benefits
varying by stratum and individual responsiveness. The intended finite-cohort
target and the standardized estimator are
\begin{equation}
\tau_E=\frac{1}{|E|}\sum_{i\in E}\{Y_i(1)-Y_i(0)\},\qquad
\hat\tau(w)=\sum_g w_g(\bar Y_{1g}-\bar Y_{0g}).
\label{eq:ecatarget}
\end{equation}
For enrolled records $E$ and controls $C$, enrolled weights are
$w_g=n_{Eg}/|E|$; pooled weights are
$w_g=(n_{Eg}+n_{Cg})/(|E|+|C|)$. Deliberately using pooled weights in both
estimator and observed-data reference makes their agreement an identity.
It is a pedagogical construction, not an empirical frequency claim.

\begin{table}[htbp]
\centering\small
\caption{Reference choice reverses the estimator ranking in the synthetic
external-control example. Bias and RMSE use 200 studies with 1{,}600 enrolled
records and 4{,}800 controls. Columns identify the population used to average
observed stratum contrasts; only enrolled standardization matches the intended
target. All entries are in abstract outcome units.}
\label{tab:eca}
\begin{tabular}{@{}lrrrr@{}}
\toprule
& \multicolumn{2}{c}{Pooled standardization} & \multicolumn{2}{c}{Enrolled standardization} \\
\cmidrule(lr){2-3}\cmidrule(l){4-5}
Validation reference & Bias & RMSE & Bias & RMSE \\
\midrule
Observed contrast, pooled & 0.000 & 0.000 & -0.872 & 0.876 \\
Potential outcomes, pooled & +0.003 & 0.115 & -0.869 & 0.875 \\
Potential outcomes, enrolled & +0.885 & 0.892 & +0.013 & 0.105 \\
\bottomrule
\end{tabular}
\end{table}


The informative control breaks that identity: the pooled potential-outcome
reference still favors pooled standardization, whereas $\tau_E$ reverses
the ranking (Table~\ref{tab:eca}). Thus avoiding an identical reference is
insufficient when it still targets the wrong population. The Monte Carlo
study quantifies errors under this chosen case mix; it does not discover
the algebraic equality. Under zero case-mix shift, pooled-standardization
bias against $\tau_E$ is only $-0.002$ units despite the same identity.
Under the full shift, weighting balances measured strata between arms but
leaves a maximum standardized difference of 0.964 against the enrolled
population. Between-arm balance does not establish population alignment.
Appendix~\ref{app:eca} specifies the 200-study-per-size design.
